\begin{abstract}
本講義旨在深入探討 Knuth-Morris-Pratt (KMP) 字串搜尋演算法。KMP 演算法是一種高效的字串匹配演算法，相較於樸素（Naive）演算法，它避免了在匹配失敗時不必要的回溯，從而顯著提高了效率。本講義將從樸素演算法的缺點切入，逐步推導 KMP 演算法的核心思想：利用模式串自身的特性（即「最長共同前後綴」）來構建一個「部分匹配表」（也稱為「失配函數」或「前綴函數」），並詳細解釋其構建過程與應用。最後，提供完整的 C++ 實作範例，幫助讀者深入理解並掌握 KMP 演算法。
\end{abstract}

\section{引言}
字串搜尋是電腦科學中一個非常基礎且重要的問題：如何在一個長字串（文本，$T$）中找到一個短字串（模式串，$P$）的所有出現位置？
最直觀的方法是樸素（Naive）演算法，但其在某些情況下效率較低。KMP 演算法正是為了解決這個效率問題而提出。

\subsection{樸素字串搜尋演算法的缺陷}
考慮在文本 $T = \text{\verb|ABABCABAB|}$ 中搜尋模式串 $P = \text{\verb|ABCABAB|}$。
\begin{itemize}
    \item 樸素演算法從 $T$ 的開頭開始，逐一比較 $P$ 和 $T$ 的子字串。
    \item 如果在某個位置不匹配，則將 $P$ 向右移動一位，重新從 $P$ 的開頭與 $T$ 的當前位置進行比較。
\end{itemize}

例如：
\begin{itemize}
    \item $T$: \verb|A B A B C A B A B|
    \item $P$: \verb|A B C A B A B|
    \item 第一次匹配：\verb|T[0]| 與 \verb|P[0]| 匹配 (\verb|A=A|)，\verb|T[1]| 與 \verb|P[1]| 匹配 (\verb|B=B|)。
    \item \verb|T[2]| (\verb|A|) 與 \verb|P[2]| (\verb|C|) 不匹配。此時，樸素演算法會將 $P$ 向右移動一位，從 \verb|T[1]| 重新開始比較。
\end{itemize}
這種「回溯」行為會導致大量的重複比較，尤其是在模式串中存在重複字元時，效率會非常低。其最差時間複雜度可達 $O(m \times n)$，其中 $n$ 是文本長度，$m$ 是模式串長度。

\section{KMP 字串搜尋演算法}
KMP 演算法的關鍵在於：當匹配失敗時，我們不需要將模式串 $P$ 完全回溯到起始位置，而是可以利用已經匹配成功的字元資訊，將 $P$ 向右「跳躍」一定的距離。這個「跳躍」的距離是由模式串 $P$ 自身的結構決定的。

KMP 演算法的核心是構建一個 \textbf{「部分匹配表」（Partial Match Table, PMT）}，也稱為 \textbf{「失配函數」（Failure Function）} 或 \textbf{「前綴函數」（Prefix Function）}。這個表紀錄了模式串 $P$ 中每個前綴的最長「共同前後綴」的長度。

\subsection{「共同前後綴」的概念}
對於一個字串 $S$，它的\textbf{前綴}是指從字串開頭到任意位置的子字串，例如 \verb|ABC| 的前綴有 \verb|A|, \verb|AB|, \verb|ABC|。它的\textbf{後綴}是指從任意位置到字串結尾的子字串，例如 \verb|ABC| 的後綴有 \verb|C|, \verb|BC|, \verb|ABC|。

\textbf{共同前後綴}：一個字串 $S$ 的一個前綴同時也是 $S$ 的一個後綴。
例如，對於字串 $S = \text{\verb|ABCAB|}$：
\begin{itemize}
    \item 前綴：\verb|A|, \verb|AB|, \verb|ABC|, \verb|ABCA|, \verb|ABCAB|
    \item 後綴：\verb|B|, \verb|AB|, \verb|CAB|, \verb|BCAB|, \verb|ABCAB|
\end{itemize}
其共同前後綴有 \verb|AB|。最長共同前後綴的長度是 2。

\subsection{部分匹配表（PMT）的定義}
\textbf{PMT[i]} 表示模式串 $P$ 的長度為 $i+1$ 的前綴 $P[0 \dots i]$ 中，最長的\textbf{真共同前後綴}（即不等於前綴本身的共同前後綴）的長度。

例如，對於模式串 $P = \text{\verb|ABCAB|}$：
\begin{itemize}
    \item \verb|P[0...0]| (\verb|A|): 最長真共同前後綴長度為 0。PMT[0] = 0。
    \item \verb|P[0...1]| (\verb|AB|): 最長真共同前後綴長度為 0。PMT[1] = 0。
    \item \verb|P[0...2]| (\verb|ABC|): 最長真共同前後綴長度為 0。PMT[2] = 0。
    \item \verb|P[0...3]| (\verb|ABCA|): 最長真共同前後綴為 \verb|A|，長度為 1。PMT[3] = 1。
    \item \verb|P[0...4]| (\verb|ABCAB|): 最長真共同前後綴為 \verb|AB|，長度為 2。PMT[4] = 2。
\end{itemize}
所以，PMT 表為 \verb|[0, 0, 0, 1, 2]|。

通常，我們使用的 PMT 陣列會根據索引調整，使得 $\text{PMT}[k]$ 對應的是長度為 $k$ 的前綴的最長共同前後綴的長度。此時 PMT 的大小為 $m$， $\text{PMT}[0]$ 沒有實際意義或定義為 -1。在 C++ 實作中，我們通常稱其為 \textbf{\verb|lps|} (longest proper prefix suffix) 陣列，且 \verb|lps[i]| 表示 $P[0 \dots i-1]$ 這個長度為 $i$ 的前綴的最長共同前後綴的長度。

例如，對於模式串 $P = \text{\verb|ABABCABAB|}$ (長度 $m=9$)
\begin{tabular}{|c|c|c|c|}
    \hline
    $i$ & $P[0 \dots i-1]$ (前綴) & 最長真共同前後綴 & `lps[i]` \\
    \hline
    0 & (空字串) & & 0 (或 -1) \\
    1 & \verb|A| & & 0 \\
    2 & \verb|AB| & & 0 \\
    3 & \verb|ABA| & \verb|A| & 1 \\
    4 & \verb|ABAB| & \verb|AB| & 2 \\
    5 & \verb|ABABC| & & 0 \\
    6 & \verb|ABABCA| & \verb|A| & 1 \\
    7 & \verb|ABCABA| & \verb|A| & 1 \\
    8 & \verb|ABCABAB| & \verb|AB| & 2 \\
    9 & \verb|ABCABABA| & \verb|ABA| & 3 \\
    \hline
\end{tabular}

實際應用中，為了簡化邏輯，\verb|lps[i]| 通常表示的是模式串 $P$ 的\textbf{長度為 $i$ 的前綴}（即 $P[0 \dots i-1]$）的最長共同前後綴的長度。對於 \verb|P[0]| (長度為 1 的前綴)，\verb|lps[1]| = 0。

\subsection{構建部分匹配表（PMT / `lps` 陣列）}
構建 \verb|lps| 陣列是 KMP 演算法的關鍵一步。這個過程本身就可以看作是一個簡化的 KMP 演算法。

我們使用兩個指標：
\begin{itemize}
    \item \verb|len|: 表示當前正在考慮的前綴的最長共同前後綴的長度。
    \item \verb|i|: 遍歷模式串的當前字元索引。
\end{itemize}

初始化：\verb|lps[0] = 0| (對於 \verb|P[0]|，沒有真共同前後綴)，\verb|len = 0|，\verb|i = 1|。

迭代過程：
對於 \verb|i| 從 1 到 $m-1$（$m$ 是模式串長度）：
\begin{enumerate}
    \item 如果 \verb|P[i] == P[len]|：
    這表示當前字元與長度為 \verb|len| 的前綴的下一個字元匹配。那麼最長共同前後綴的長度可以增加。
    \verb|len++|；\verb|lps[i] = len|；\verb|i++|。
    \item 如果 \verb|P[i] != P[len]|：
    這表示不匹配。我們需要尋找一個更短的共同前後綴。
    \begin{itemize}
        \item 如果 \verb|len != 0|：
            我們不能簡單地將 \verb|len| 設為 0，而是需要根據 \verb|lps[len-1]| 的值，回溯到前一個最長共同前後綴的長度。這樣做是因為 \verb|lps[len-1]| 儲存了 \verb|P[0...len-1]| 這個前綴的最長共同前後綴的長度，相當於在 \verb|P[0...len-1]| 中尋找下一個可能的匹配起點。
            \verb|len = lps[len-1]|。
        \item 如果 \verb|len == 0|：
            這表示已經回溯到最開頭，沒有更短的共同前後綴了。那麼 \verb|P[0...i]| 的最長共同前後綴長度就是 0。
            \verb|lps[i] = 0|；\verb|i++|。
    \end{itemize}
\end{enumerate}

\subsubsection{構建 `lps` 陣列的範例}
模式串 $P = \text{\verb|AAACAAAAAC|}$ ($m=10$)
初始化：\verb|lps| 陣列大小為 10，\verb|len = 0|，\verb|i = 1|。

\begin{enumerate}
    \item \verb|i = 1|, \verb|P[1] = 'A'|, \verb|P[len] = P[0] = 'A'|
    \verb|P[1] == P[0]|。
    \verb|len| 變為 1。\verb|lps[1] = 1|。\verb|i| 變為 2。
    \verb|lps = [0, 1, ?, ?, ?, ?, ?, ?, ?, ?]|

    \item \verb|i = 2|, \verb|P[2] = 'A'|, \verb|P[len] = P[1] = 'A'|
    \verb|P[2] == P[1]|。
    \verb|len| 變為 2。\verb|lps[2] = 2|。\verb|i| 變為 3。
    \verb|lps = [0, 1, 2, ?, ?, ?, ?, ?, ?, ?]|

    \item \verb|i = 3|, \verb|P[3] = 'C'|, \verb|P[len] = P[2] = 'A'|
    \verb|P[3] != P[2]|。
    \verb|len != 0| (len=2)。所以 \verb|len = lps[len-1] = lps[1] = 1|。
    \verb|lps = [0, 1, 2, ?, ?, ?, ?, ?, ?, ?]| (len 變回 1)

    \item \verb|i = 3|, \verb|P[3] = 'C'|, \verb|P[len] = P[1] = 'A'|
    \verb|P[3] != P[1]|。
    \verb|len != 0| (len=1)。所以 \verb|len = lps[len-1] = lps[0] = 0|。
    \verb|lps = [0, 1, 2, ?, ?, ?, ?, ?, ?, ?]| (len 變回 0)

    \item \verb|i = 3|, \verb|P[3] = 'C'|, \verb|P[len] = P[0]| (len=0，這表示比較 \verb|P[3]| 和 \verb|P[0]|)
    \verb|P[3] != P[0]|。
    \verb|len == 0|。
    \verb|lps[3] = 0|。\verb|i| 變為 4。
    \verb|lps = [0, 1, 2, 0, ?, ?, ?, ?, ?, ?]|

    \item \verb|i = 4|, \verb|P[4] = 'A'|, \verb|P[len] = P[0] = 'A'|
    \verb|P[4] == P[0]|。
    \verb|len| 變為 1。\verb|lps[4] = 1|。\verb|i| 變為 5。
    \verb|lps = [0, 1, 2, 0, 1, ?, ?, ?, ?, ?]|

    \item \verb|i = 5|, \verb|P[5] = 'A'|, \verb|P[len] = P[1] = 'A'|
    \verb|P[5] == P[1]|。
    \verb|len| 變為 2。\verb|lps[5] = 2|。\verb|i| 變為 6。
    \verb|lps = [0, 1, 2, 0, 1, 2, ?, ?, ?, ?]|

    \item \verb|i = 6|, \verb|P[6] = 'A'|, \verb|P[len] = P[2] = 'A'|
    \verb|P[6] == P[2]|。
    \verb|len| 變為 3。\verb|lps[6] = 3|。\verb|i| 變為 7。
    \verb|lps = [0, 1, 2, 0, 1, 2, 3, ?, ?, ?]|

    \item \verb|i = 7|, \verb|P[7] = 'A'|, \verb|P[len] = P[3] = 'C'|
    \verb|P[7] != P[3]|。
    \verb|len != 0| (len=3)。所以 \verb|len = lps[len-1] = lps[2] = 2|。
    \verb|lps = [0, 1, 2, 0, 1, 2, 3, ?, ?, ?]| (len 變回 2)

    \item \verb|i = 7|, \verb|P[7] = 'A'|, \verb|P[len] = P[2] = 'A'|
    \verb|P[7] == P[2]|。
    \verb|len| 變為 3。\verb|lps[7] = 3|。\verb|i| 變為 8。
    \verb|lps = [0, 1, 2, 0, 1, 2, 3, 3, ?, ?]|

    \item \verb|i = 8|, \verb|P[8] = 'A'|, \verb|P[len] = P[3] = 'C'|
    \verb|P[8] != P[3]|。
    \verb|len != 0| (len=3)。所以 \verb|len = lps[len-1] = lps[2] = 2|。
    \verb|lps = [0, 1, 2, 0, 1, 2, 3, 3, ?, ?]| (len 變回 2)

    \item \verb|i = 8|, \verb|P[8] = 'A'|, \verb|P[len] = P[2] = 'A'|
    \verb|P[8] == P[2]|。
    \verb|len| 變為 3。\verb|lps[8] = 3|。\verb|i| 變為 9。
    \verb|lps = [0, 1, 2, 0, 1, 2, 3, 3, 3, ?]|

    \item \verb|i = 9|, \verb|P[9] = 'C'|, \verb|P[len] = P[3] = 'C'|
    \verb|P[9] == P[3]|。
    \verb|len| 變為 4。\verb|lps[9] = 4|。\verb|i| 變為 10 (結束)。
    \verb|lps = [0, 1, 2, 0, 1, 2, 3, 3, 3, 4]|
\end{enumerate}

最終 \verb|lps| 陣列為 \verb|[0, 1, 2, 0, 1, 2, 3, 3, 3, 4]|。

\subsection{KMP 演算法匹配過程}
有了 \verb|lps| 陣列後，KMP 的匹配過程就變得高效。

我們使用兩個指標：
\begin{itemize}
    \item \verb|i|: 指向文本串 $T$ 的當前字元索引。
    \item \verb|j|: 指向模式串 $P$ 的當前字元索引。
\end{itemize}

初始化：\verb|i = 0|，\verb|j = 0|。

迭代過程：
當 \verb|i < n| 且 \verb|j < m| ($n$ 是文本長度，$m$ 是模式串長度)：
\begin{enumerate}
    \item 如果 \verb|P[j] == T[i]|：
    這表示匹配成功。繼續比較下一個字元。
    \verb|i++|；\verb|j++|。
    \item 如果 \verb|P[j] != T[i]|：
    這表示匹配失敗。
    \begin{itemize}
        \item 如果 \verb|j != 0|：
            我們不需要回溯文本串 $T$ 的 \verb|i|，而是利用 \verb|lps| 陣列將模式串 $P$ 向右「跳躍」。新的 \verb|j| 值將是 \verb|lps[j-1]|。這表示我們將模式串的 \verb|P[0...j-1]| 這個前綴的最長共同前後綴對齊到 $T$ 的當前匹配位置。
            \verb|j = lps[j-1]|。
        \item 如果 \verb|j == 0|：
            這表示模式串的第一個字元就沒有匹配成功。我們只能將文本串 $T$ 的指針向右移動一位，模式串 $P$ 的指針依然保持在 0。
            \verb|i++|。
    \end{itemize}
\end{enumerate}

匹配成功判斷：
如果 \verb|j == m|，表示模式串 $P$ 已經在文本串 $T$ 中找到了一個匹配。此時匹配的起始位置是 \verb|i - m|。
找到一個匹配後，為了尋找下一個匹配，我們需要將 \verb|j| 重置為 \verb|lps[j-1]|，繼續從當前 $T$ 的位置尋找。

\subsubsection{KMP 匹配範例}
文本 $T = \text{\verb|ABABCABABABABABCABAB|}$
模式串 $P = \text{\verb|ABCABAB|}$
其 \verb|lps| 陣列為 \verb|[0, 0, 0, 1, 2, 3, 4]|

初始化：\verb|i = 0|, \verb|j = 0|

\begin{enumerate}
    \item \verb|T[0]='A'|, \verb|P[0]='A'| $\rightarrow$ 匹配。\verb|i=1, j=1|
    \item \verb|T[1]='B'|, \verb|P[1]='B'| $\rightarrow$ 匹配。\verb|i=2, j=2|
    \item \verb|T[2]='A'|, \verb|P[2]='C'| $\rightarrow$ 不匹配。
    \verb|j != 0| (j=2)。\verb|j = lps[j-1] = lps[1] = 0|。
    此時 $T$ 指針 \verb|i| 不變，仍為 2。模式串 $P$ 移位，從 \verb|P[0]| 重新與 \verb|T[2]| 比較。
    \verb|i=2, j=0|

    \item \verb|T[2]='A'|, \verb|P[0]='A'| $\rightarrow$ 匹配。\verb|i=3, j=1|
    \item \verb|T[3]='B'|, \verb|P[1]='B'| $\rightarrow$ 匹配。\verb|i=4, j=2|
    \item \verb|T[4]='C'|, \verb|P[2]='C'| $\rightarrow$ 匹配。\verb|i=5, j=3|
    \item \verb|T[5]='A'|, \verb|P[3]='A'| $\rightarrow$ 匹配。\verb|i=6, j=4|
    \item \verb|T[6]='B'|, \verb|P[4]='B'| $\rightarrow$ 匹配。\verb|i=7, j=5|
    \item \verb|T[7]='A'|, \verb|P[5]='A'| $\rightarrow$ 匹配。\verb|i=8, j=6|
    \item \verb|T[8]='B'|, \verb|P[6]='B'| $\rightarrow$ 匹配。\verb|i=9, j=7|
    此時 \verb|j == m| (j=7, m=7)。找到一個匹配！
    匹配起始位置為 \verb|i - m = 9 - 7 = 2|。
    為了尋找下一個匹配，\verb|j = lps[j-1] = lps[6] = 4|。
    \verb|i=9, j=4| (繼續比較 \verb|T[9]| 和 \verb|P[4]|)

    \item \verb|T[9]='A'|, \verb|P[4]='B'| $\rightarrow$ 不匹配。
    \verb|j != 0| (j=4)。\verb|j = lps[j-1] = lps[3] = 1|。
    \verb|i=9, j=1|

    \item \verb|T[9]='A'|, \verb|P[1]='B'| $\rightarrow$ 不匹配。
    \verb|j != 0| (j=1)。\verb|j = lps[j-1] = lps[0] = 0|。
    \verb|i=9, j=0|

    \item \verb|T[9]='A'|, \verb|P[0]='A'| $\rightarrow$ 匹配。\verb|i=10, j=1|
    ... 繼續匹配直到文本結束。
\end{enumerate}

\subsection{時間複雜度分析}
\begin{itemize}
    \item \textbf{構建 \verb|lps| 陣列：}
        兩個指針 \verb|i| 和 \verb|len| 都單向遞增，並且 \verb|len| 在每次不匹配時會通過 \verb|lps[len-1]| 縮短，但總的縮短次數不會超過 \verb|len| 的總增長次數。因此，構建 \verb|lps| 陣列的時間複雜度是 $O(m)$，其中 $m$ 是模式串的長度。
    \item \textbf{KMP 匹配過程：}
        類似地，文本指針 \verb|i| 和模式串指針 \verb|j| 都單向遞增。當不匹配時，\verb|j| 會通過 \verb|lps[j-1]| 縮短，但總的縮短次數不會超過 \verb|j| 的總增長次數。因此，KMP 匹配過程的時間複雜度是 $O(n)$，其中 $n$ 是文本的長度。
    \item \textbf{總體時間複雜度：} $O(m + n)$。
    \item \textbf{空間複雜度：} $O(m)$，用於儲存 \verb|lps| 陣列。
\end{itemize}
相較於樸素演算法的 $O(m \times n)$ 最差時間複雜度，KMP 演算法在效率上有了顯著提升。

\subsection{C++ 實作}

\lstset{caption={KMP 演算法的 C++ 實作}}
\begin{lstlisting}[language=C++]
#include <bits/stdc++.h>
using namespace std;

vector<int> computeLPSArray(const string& pat) {
    int m = pat.length();
    vector<int> lps(m, 0);
    int length = 0;
    int i = 1;

    // Loop calculates lps[i] for i = 1 to m-1
    while (i < m) {
        if (pat[i] == pat[length]) {
            length++;
            lps[i] = length;
            i++;
        } else { // (pat[i] != pat[length])
            if (length != 0) {
                length = lps[length - 1];
            } else {
                lps[i] = 0;
                i++;
            }
        }
    }
    return lps;
}

void KMPSearch(const string& txt, const string& pat) {
    int n = txt.length();
    int m = pat.length();

    if (m == 0) {
        cout << "Pattern is empty. Found at every position." << endl;
        return;
    }
    if (n == 0) {
        cout << "Text is empty. No pattern found." << endl;
        return;
    }
    if (m > n) {
        cout << "Pattern length is greater than text length. No match possible." << endl;
        return;
    }

    vector<int> lps = computeLPSArray(pat);

    cout << "LPS array for pattern \"" << pat << "\": [";
    for (int k = 0; k < m; ++k) {
        cout << lps[k] << (k == m - 1 ? "" : ", ");
    }
    cout << "]" << endl;

    int i = 0, j = 0;
    int count = 0;

    while (i < n) {
        if (pat[j] == txt[i]) {
            i++;
            j++;
        }

        if (j == m) {
            cout << "Found pattern at index " << i - j << endl;
            count++;
            j = lps[j - 1];
        } else if (i < n && pat[j] != txt[i]) {
            if (j != 0) {
                j = lps[j - 1];
            } else {
                i++;
            }
        }
    }

    if (count == 0) {
        cout << "No pattern found in the text." << endl;
    }
}

int main() {
    string text1 = "ABABDABACDABABCABAB";
    string pattern1 = "ABABCABAB";
    cout << "Text: \"" << text1 << "\"" << endl;
    cout << "Pattern: \"" << pattern1 << "\"" << endl;
    KMPSearch(text1, pattern1);
    cout << "--------------------" << endl;

    string text2 = "AAAAABAAABAA";
    string pattern2 = "AAAB";
    cout << "Text: \"" << text2 << "\"" << endl;
    cout << "Pattern: \"" << pattern2 << "\"" << endl;
    KMPSearch(text2, pattern2);
    cout << "--------------------" << endl;

    string text3 = "ABCDEFGH";
    string pattern3 = "XYZ";
    cout << "Text: \"" << text3 << "\"" << endl;
    cout << "Pattern: \"" << pattern3 << "\"" << endl;
    KMPSearch(text3, pattern3);
    cout << "--------------------" << endl;

    string text4 = "ABABABABC";
    string pattern4 = "ABC";
    cout << "Text: \"" << text4 << "\"" << endl;
    cout << "Pattern: \"" << pattern4 << "\"" << endl;
    KMPSearch(text4, pattern4);
    cout << "--------------------" << endl;

    return 0;
}
\end{lstlisting}

\section{總結}
KMP 演算法利用模式串自身的重複特性，避免了樸素演算法中不必要的回溯，從而將字串搜尋的時間複雜度優化到線性的 $O(m+n)$。其核心在於構建一個「部分匹配表」（LPS 陣列），該表指導了模式串在不匹配時如何高效地跳躍。掌握 KMP 演算法不僅能解決字串匹配問題，其背後的動態規劃思想和對問題結構的利用，對於理解其他演算法也大有裨益。
